
\documentclass[12pt]{article}
\pagestyle{empty} \topmargin -0.5in \textheight 9in \textwidth 6.5in
\oddsidemargin -0in
\usepackage{amssymb}
\usepackage{amsmath}
\usepackage[normalem]{ulem}
\usepackage{enumerate}
\usepackage{xcolor}
\DeclareMathOperator{\tr}{tr}

\newcommand{\vect}[1]{\mathbf{#1}}

\begin{document}

\begin{flushleft}

\textbf{Department of Applied Mathematics and Statistics}

\textbf{COLORADO SCHOOL OF MINES}\\


MATH $455$: Partial Differential Equations
\end{flushleft}

\vspace{0.1in}

\begin{center}
{\bf Assignment \#5 - Laplace's Equation} \\
Due Thursday, November 5, 2026 \\
\end{center}

\vspace{0.3in}

\noindent $\textbf{1}$. (14 points) Let $a,b > 0$ and $g \in \mathrm{PC}^1[0,a]$ be given. Find the solution $u(x,y)$ of 
$$\Delta u = 0$$
on the rectangle $ \Omega = [0,a] \times [0,b],$
with
$$\frac{\partial u}{\partial x}(0,y) = \frac{\partial u}{\partial x} (a,y) = 0$$
for $0 < y < b$ and
$$\frac{\partial u}{\partial y}(x,0) = 0 \qquad \mathrm{and} \qquad u(x,b) = g(x)$$
for $0 < x < a$.

\noindent \underline{\emph{Hint:}} It might be easiest to use $\sinh$ and $\cosh$ in some places rather than real exponentials. \\ \\
\textcolor{blue}{\textbf{Solution} — 
\begin{align*}
\Delta u = 0 \implies& \frac{\partial^2 u}{\partial x^2} + \frac{\partial^2 u}{\partial y^2} = 0
\end{align*}
Since this is a linear PDE, we proceed with separation of variables, assuming the solution takes the (separated) form $u(x,y) = X(x)Y(y)$ \\
\begin{align*}
\implies& X''(x)Y(y) = -X(x)Y''(y) \\
\implies& \frac{X''(x)}{X(x)} = -\frac{Y''(y)}{Y(y)} = \lambda
\end{align*} \\
This equation takes the form $F(x) = G(y)$ for all $x,y$ $\implies$ both must be equal to some constant $\lambda$ (fix $x = x_0$, and then $F(x_0) = G(y)$ for all $y \implies G(y)$ is constant $\implies F(x)$ is constant since they're equal to each other). This yields the two ordinary differential equations
\begin{align*}
\begin{cases}
X'' - \lambda X = 0 \\
Y'' + \lambda Y = 0
\end{cases}
\end{align*} \\
whose solutions depend on the sign of $\lambda$. \\ \\
\textbf{For $\lambda < 0$}: \\ \\
$$X(x) = c_0 \cos \left( \sqrt{-\lambda} x \right) + c_1 \sin \left( \sqrt{-\lambda} x \right) \;\;\;\;\;\;\;\;\; Y(y) = c_2 e^{\sqrt{-\lambda} y} + c_3 e^{-\sqrt{-\lambda} y}$$
\textbf{For $\lambda = 0$}: \\ \\
$$X(x) = Ax + B \;\;\;\;\;\;\;\;\; Y(y) = Cy + D$$
\textbf{For $\lambda > 0$}: \\ \\
$$X(x) = c_4 e^{\sqrt{\lambda} x} + c_5 e^{-\sqrt{\lambda} x} \;\;\;\;\;\;\;\;\; Y(y) = c_6 \cos \left( \sqrt{\lambda} y \right) + c_7 \sin \left( \sqrt{\lambda} y \right)$$ \\
We apply the boundary conditions to $u(x,y) = X(x)Y(y)$: \\ \\
\begin{align*}
\frac{\partial u}{\partial x}(0,y) = 0 \implies& X'(0)Y(y) = 0 \\
\implies& X'(0) = 0 \;\;\;\; \text{(For $Y(y) \neq 0$)}
\end{align*}
\begin{align*}
\frac{\partial u}{\partial x}(a,y) = 0 \implies& X'(a)Y(y) = 0 \\
\implies& X'(a) = 0 \;\;\;\; \text{(For $Y(y) \neq 0$)}
\end{align*}
We use these results to constrain the unknown constants in our separated solutions: \\ \\
\textbf{For $\lambda < 0$}: \\ \\
\begin{align*}
X'(0) = 0 \implies& \sqrt{-\lambda} c_1 = 0 \\
\implies& \boxed{c_1 = 0} \\
X'(a) = 0 \implies& -\sqrt{-\lambda} c_0 \sin \left( \sqrt{-\lambda} a \right) = 0 \\
\implies& \sin \left( \sqrt{-\lambda} a \right) = 0 \;\;\;\; \text{(for nontrivial solution with $c_0 \neq 0$)} \\
\implies& \sqrt{-\lambda} a = n\pi \;\;\;\;\; \text{for $n \in \mathbb{Z}$} \\
\implies& \boxed{\lambda = - \left( \frac{n \pi}{a} \right) ^2}
\end{align*}
\textbf{For $\lambda = 0$}: \\
\begin{align*}
X'(0) = 0, X'(a) = 0 \implies& \boxed{A = 0}\\
\end{align*}
\textbf{For $\lambda > 0$}: \\ \\
\begin{align*}
X'(0) = 0 \implies& \sqrt{\lambda} c_4 - \sqrt{\lambda} c_5 = 0 \\
\implies& \boxed{c_4 = c_5} \\
X'(a) = 0 \implies& \sqrt{\lambda} c_4 e^{\sqrt{\lambda} a} - \sqrt{\lambda} c_5 e^{-\sqrt{\lambda} a} = 0 \\
\implies& \sqrt{\lambda} c_4 \left( \underbrace{ e^{\sqrt{\lambda} a} - e^{-\sqrt{\lambda} a} }_{\neq 0} \right) = 0 \;\;\; \text{(using $c_4 = c_5$)} \\
\implies& \boxed{c_4 = 0}
\end{align*}
This leads us to general separated solutions: \\ \\
\textbf{For $\lambda < 0$}: \\ \\
$$u(x,y) = \cos \left( \frac{n \pi x}{a} \right) \left( A_n e^{ \frac{n\pi y}{a} } + B_n e^{-\frac{n\pi y}{a}}  \right)$$
\textbf{For $\lambda = 0$}: \\ \\
$$u(x,y) = \alpha y + \beta$$
\textbf{For $\lambda > 0$}: \\ \\
$$u(x,y) = 0$$
Hence the general vertically-constrained solution is the superposition of all these separated solutions: \\
$$u(x,y) = \alpha y + \beta + \sum_{n=1}^{\infty} \cos \left( \frac{n \pi x}{a} \right) \left( A_n e^{ \frac{n\pi y}{a} } + B_n e^{-\frac{n\pi y}{a}}  \right)$$ \\
We proceed to apply the horizontal boundary conditions. \\
\begin{align*}
\frac{\partial u}{\partial y}(x,0) = 0 \implies&  \left. \alpha + \sum_{n=1}^{\infty} \frac{n\pi}{a} \cos \left( \frac{n \pi x}{a} \right) \left( A_n e^{ \frac{n\pi y}{a} } - B_n e^{-\frac{n\pi y}{a}} \right) \right|_{(x,0)} = 0 \\
\implies& \alpha + \sum_{n=1}^{\infty} \frac{n\pi}{a} (A_n - B_n) \cos \left( \frac{n \pi x}{a} \right) = 0  \\
\end{align*}
Hence $\alpha$ is the constant term and each $n \pi (A_n - B_n)  / a$ is the Fourier coefficient of the cosine expansion of $f(x) = 0$ on $x\in[0,a]$. Since the average value of this function is $0$ everyone, we have that $\boxed{\alpha = 0}$. Also, \\
\begin{align*}
\implies& \frac{n\pi (A_n - B_n) }{a} = \frac{2}{a} \int_0^a (-\alpha) \cos \left( \frac{n\pi x}{a} \right) dx \\
\implies& (A_n - B_n)  = -\frac{2\alpha}{n\pi} \int_0^a \cos \left( \frac{n\pi x}{a} \right) dx = 0 \\
\implies & \boxed{A_n = B_n}
\end{align*} \\
Hence we can write the solution as
$$u(x,y) =  \beta + 2 \sum_{n=1}^{\infty} C_n \cos \left( \frac{n \pi x}{a} \right) \cosh \left( \frac{n \pi y}{a} \right)$$ \\
We apply the final boundary condition
\begin{align*}
u(x,b) = g(x) \implies& \beta + 2 \sum_{n=1}^{\infty} C_n \cos \left( \frac{n \pi x}{a} \right) \cosh \left( \frac{n \pi b}{a} \right) = g(x) \\
\implies& g(x) = \beta + \sum_{n=1}^{\infty} \left[ \underbrace{ 2C_n \cosh \left( \frac{n \pi b}{a} \right) }_{\zeta_n} \right] \cos \left( \frac{n \pi x}{a} \right) \\
\implies& g(x) = \beta + \sum_{n=1}^{\infty} \zeta_n \cos \left( \frac{n \pi x}{a} \right)
\end{align*}
Which is the cosine-expansion of $g(x)$ on $[0,a]$, with $\beta$ being the average value of $g(x)$ on $[0,a]$. Hence
\begin{align*}
\boxed{\beta = \frac{1}{a} \int_0^a g(x) dx} \\ \\
\end{align*}
And
\begin{align*}
\zeta_n &= \frac{2}{a} \int_0^a g(x) \cos \left( \frac{n \pi x}{a} \right) dx \\
\implies& 2C_n \cosh \left( \frac{n \pi b}{a} \right) = \frac{2}{a} \int_0^a g(x) \cos \left( \frac{n \pi x}{a} \right) dx \\
\implies& \boxed{C_n = \frac{1}{a} \operatorname{sech} \left( \frac{n \pi b}{a} \right) \int_0^a g(x) \cos \left( \frac{n \pi x}{a} \right)} dx 
\end{align*}
In the general solution
$$\boxed{u(x,y) =  \beta + 2 \sum_{n=1}^{\infty} C_n \cos \left( \frac{n \pi x}{a} \right) \cosh \left( \frac{n \pi y}{a} \right)}$$
}
\vspace{0.3in}

\noindent $\textbf{2}$. (16 points) Let $a > 0$ be given and find (using polar coordinates) the solution $u(r,\theta)$ of Laplace's equation
$$\Delta u = 0$$
on the wedge
$$ \Omega = \left \{ (r,\theta) : 0 \leq r < a, 0 \leq \theta \leq \frac{\pi}{3} \right \}$$
with
$$u(r,0) = u\left (r,\frac{\pi}{3} \right)  = 0$$
for $0 \leq r < a$ and
$$u(a,\theta) = a^2$$
for $0 \leq \theta \leq \frac{\pi}{3}$.
In particular, make sure to only find solutions that remain bounded as $r \to 0$.

\vspace{0.3in}


\noindent $\textbf{3}$. (20 points) Let $0 < R_1 < R_2$ be given and $g$ be a PC$^1$ function. Find the single-valued solution of Laplace's equation on the annulus
$$ \Omega = \left \{ (r,\theta) : R_1 < r < R_2, -\pi \leq \theta < \pi \right \}$$
with the boundary conditions
$$u(R_1, \theta) = 0 \qquad \mathrm{and} \qquad u(R_2, \theta) = g(\theta)$$
for $-\pi \leq \theta < \pi$.
In particular, use the periodic boundary conditions
$$u(r, -\pi) = u(r, \pi) \qquad \mathrm{and} \qquad \frac{\partial u}{\partial \theta}(r, -\pi) = \frac{\partial u}{\partial \theta}(r, \pi)$$\
for $R_1 < r < R_2$
in order to compute the solution.


\end{document}
